\section{Evaluation and Results}

\subsection{Evidence Model}

The evaluation separates five evidence layers: source conformance, deterministic rule behaviour, adversarial boundary behaviour, Android package integration, and physical-device usability evidence. This separation prevents a passed unit test from being reported as device reliability or legal validity. The final acceptance record is stored at \path{testing-report/release-candidate-2026-08-09.md}; screenshots and UI hierarchies are retained in the adjacent dated folder.

\subsection{Source and Rule Verification}

The complete application passed \texttt{tsc --noEmit} and ESLint. The compliance sources were compiled before execution, preventing stale generated JavaScript from being treated as current evidence.

A fixed-seed logical campaign generated 1,000 cases. Eight hundred were expected review signals and 200 were below-threshold controls. The observed confusion matrix was

\begin{equation}
(TP,TN,FP,FN)=(800,200,0,0),
\end{equation}

which gives precision and recall of 1.0000 against the specified oracle. The oracle and engine implement the same published threshold model; this is a conformance and regression result, not an independent estimate of real-world infringement detection. Zero observed errors do not establish a zero population error rate.

Extended suites passed for unsafe types and integers, malformed context, unknown source, permission-key mutation, temporal overlap, adjacency, replay, retention, simulator isolation, dashboard presentation, missing-context disclosure, and regulation-pack identity. The tests confirm the enumerated boundaries. They cannot establish correctness for telemetry the application never receives or for legal interpretation outside the fixtures.

\subsection{Release Build and Manifest Verification}

Gradle completed \texttt{assembleRelease} and \texttt{bundleRelease} with target SDK 36. Table~\ref{tab:r9-artifacts} records the generated artefacts.

\begin{table}[H]
\centering
\small
\caption{Revision 9 Android release artefacts}
\label{tab:r9-artifacts}
\begin{tabular}{p{0.31\textwidth}r p{0.42\textwidth}}
\toprule
Artefact & Bytes & SHA-256 \\
\midrule
\texttt{app-release.apk} & 62,762,819 & \texttt{\seqsplit{4A5BCED7FBF6A7A5A0E337BFAADB003A3C7FE89FDAFBF9BADCBE6BC4CBFF51B5}} \\
\texttt{app-release.aab} & 31,700,025 & \texttt{\seqsplit{E4090FEEF4EFBD0B697F863D80D28F289C3554BA89394A1435EACA2C226ACF78}} \\
\bottomrule
\end{tabular}
\end{table}

The merged release manifest contained declarations for Internet access, wake lock, network state, boot rescheduling, foreground service, and an app-scoped signature permission. External-storage read and write permissions were absent. Inspection of the installed package confirmed the same requested-permission set and no runtime permissions. Android backup was disabled in the application declaration.

\subsection{Physical-Device Evaluation}

The release APK was installed by ADB on one OPPO PERM00 ARM64 device. The installed package reported \texttt{com.zhihengzhang.privacylens}, versionName 1.1.0, versionCode 2, minSdk 24, and targetSdk 36. Launcher cold start succeeded. A filtered scan of recent \texttt{AndroidRuntime} and \texttt{ReactNativeJS} error channels returned no fatal match during the acceptance run.

The three primary destinations rendered at the device's captured application resolution. Coordinate-driven exploration initially appeared to show a navigation-target problem because the physical display and captured application coordinate spaces differed. A bounded coordinate sweep and UI-tree inspection identified the offset; subsequent navigation and custom bottom-bar interaction succeeded. This calibration episode is retained because it demonstrates why raw ADB coordinates must not be interpreted without the target's current window geometry.

The controlled demonstration completed a 50-record synthetic workflow. The populated overview displayed four review items, zero evidence gaps, two no-concern items, and descriptive precision and recall of 100\% for the simulator's known labels. The Findings screen marked the source as synthetic and displayed the EU GDPR pack. These values validate the demonstration path only; they are not device-observation accuracy.

\subsection{Usability and Trust Assessment}

The product was assessed against three release-candidate questions.

\begin{table}[H]
\centering
\small
\caption{Bounded product assessment}
\begin{tabular}{>{\raggedright\arraybackslash}p{0.19\textwidth}>{\raggedright\arraybackslash}p{0.32\textwidth}>{\raggedright\arraybackslash}p{0.37\textwidth}}
\toprule
Criterion & Supporting evidence & Boundary \\
\midrule
Visually coherent & consistent brand, hierarchy, spacing, cards, iconography, and three-destination navigation on the tested device & one device; no participant preference study \\
Trustworthy & local-first disclosure, source and pack labels, explicit missing evidence, no legal-verdict language & not a legal audit or security certification \\
Willing to use & clear first action, concise summaries, optional details, demonstration, and direct data clearing & evaluator judgement; no longitudinal adoption measure \\
\bottomrule
\end{tabular}
\end{table}

The evaluator judged the candidate to meet the project's ``beautiful, trustworthy, and willing to use'' gate for an initial submission candidate. This is a structured expert judgement supported by rendered evidence, not a statistically generalisable user study.

\subsection{Evaluation Questions and Acceptance Logic}

The evaluation was organised around six questions. First, does the source compile and satisfy static rules? Second, does the decision engine agree with its specified thresholds and statuses? Third, does it reject malformed or adversarial evidence safely? Fourth, do regulation packs remain independent and traceable? Fifth, can the Android release artefacts be assembled, inspected, installed, and launched? Sixth, does the rendered product communicate the expected hierarchy and evidence boundaries on the available physical device?

Each question has a distinct acceptance condition. Source acceptance requires a clean TypeScript application check and ESLint run. Rule acceptance requires exact expected confusion-matrix values for the fixed corpus and passing boundary assertions. Runtime-boundary acceptance requires every malformed case to return the specified rejection code without an uncaught exception. Pack acceptance requires registry allowlisting, different thresholds where defined, retained pack identity, and a non-legal caveat for the Research Baseline.

Android build acceptance requires both release tasks to complete and the artefacts to exist with recorded hashes. Permission acceptance requires agreement between the merged manifest and installed package, with legacy external-storage permissions absent. Physical acceptance requires successful streamed installation, launcher start, stable primary navigation, captured final states, and no matching fatal error in the bounded log scan. Visual acceptance requires inspection of actual screenshots rather than source code alone.

The protocol deliberately does not set an acceptance condition for GDPR compliance, legal accuracy, population detection accuracy, long-run scheduler reliability, battery consumption, multi-device compatibility, or store approval. Those outcomes require evidence not collected in this campaign.

\subsection{Reproducibility Controls}

The compliance sources are compiled immediately before the generated JavaScript runners execute. This control was introduced because a failed TypeScript command can otherwise leave older generated output in place, producing a misleading green test. The acceptance sequence checks the compiler exit code before invoking Node. Generated compliance output is ignored by Git so it cannot become an accidental source of truth.

The main logical corpus uses a fixed linear-congruential seed. A fixed seed makes a failure reproducible while still exploring varied generated values. The 1,000 rounds preserve an 80/20 design: four of every five cases are expected review signals and one is a control. A fresh engine is created for each main-corpus case so temporal state from one generated subject cannot leak into another.

Boundary cases are not left to random chance. For every supported permission, counts at threshold minus two, minus one, equal, plus one, and plus two are exercised. The assertion encodes the strict boundary: only values above the prompt threshold become active under the daily-total rule. Burst tests similarly check exactly at and exactly above each permission-specific limit.

Time is fixed in extended tests. Windows and timestamp arrays are created relative to a constant epoch value rather than the wall clock. This prevents a test from becoming invalid while it runs and makes future failures comparable. Production validation still compares implausible future times against the current clock, but fixtures choose values far enough in the past to avoid accidental rejection.

\subsection{Deterministic Corpus Interpretation}

The confusion matrix compares the engine's active-state decision with the generator's specified expected label:

\begin{equation}
\mathrm{Precision}=\frac{TP}{TP+FP}, \qquad
\mathrm{Recall}=\frac{TP}{TP+FN}.
\end{equation}

The resulting precision and recall of 1.0000 mean that the implementation agreed with the oracle for all generated cases. This supports a regression claim: later refactoring did not alter the expected threshold semantics for the corpus. It also supports a fixture-quality claim: the corpus contains both positive and negative cases and the metric implementation returns the expected counts.

It does not show that a threshold crossing identifies unlawful processing. The expected label is generated from the same documented research model that the engine is meant to implement. Although the oracle is expressed independently of the evaluated result, both are derived from one specification. A defect in that specification could be reproduced consistently. There is no independently labelled sample of real applications or legal cases.

The tests also verify metric edge conditions. Empty samples, all-negative samples, and samples with missed positives must produce finite precision and recall rather than not-a-number or infinity. This matters because a dashboard can otherwise display an apparently precise but meaningless metric when a denominator is zero.

\subsection{Adversarial Input Results}

The initial security cases reject unsupported permission values, not-a-number counts, negative counts, reversed windows, and prototype-like permission keys. Extended cases add unsafe integers and positive infinity. Every such count returns the invalid-count code. Timestamp evidence is rejected when array length differs from count, when an element falls before the start, when it falls after the end, or when it is fractional.

Processing context is attacked through type confusion. Numeric purpose, array controller identity, object Article 9 condition, string consent-withdrawn, numeric special-category, and string user-initiated values all return invalid context. An unknown lawful basis and negative retention period are also rejected. These cases are significant because JavaScript truthiness or an early string operation could otherwise convert malformed evidence into a plausible status or throw outside the controlled error path.

Source validation rejects a forged source. This prevents a caller from using an arbitrary label that resembles an authoritative acquisition path. Package and permission history are isolated: a high-count record for one package cannot make another package active. Histories are also isolated by permission family.

The assertions demonstrate the implemented rejection surface, not complete security. They do not cover a compromised native module, a modified application binary, corrupted AsyncStorage beyond the migration parser, malicious operating system, or every possible Unicode identifier. Security conclusions remain bounded to the enumerated input contract.

\subsection{Temporal and Replay Results}

Burst tests create timestamp arrays exactly at and just above the location, microphone, and contacts limits. At-limit arrays produce no burst signal; one additional event produces the signal. A separate imported-evidence example concentrates 20 location events within one minute and is detected.

Cross-window tests use two below-threshold location summaries. The first remains inactive. When the second begins after the first has completed, their rolling total exceeds the threshold and the engine emits a cross-window signal. The same construction is repeated with adjacent windows to establish that equality between a prior end and current start is admissible.

An overlapping construction uses one interval that intersects the next. The second finding's rolling count remains its own count and no cross-window signal appears. This is the expected conservative behaviour because aggregate evidence cannot identify duplicate events in the overlap. An exact replay is then evaluated twice before a compatible next interval. Rolling total confirms that the repeated window contributed once, not twice.

Simulator histories are excluded from cross-window signalling. The demonstration may still produce a daily or burst signal for a record, but it cannot present accumulated simulator activity as imported temporal evidence. This is a provenance safeguard, not a statistical assumption.

The temporal tests operate in process memory and do not demonstrate persistence of aggregate history across process death. They establish the logic of the current engine instance. A production design requiring durable cross-window evidence would need a versioned evidence-history repository and migration tests separate from the visible finding repository.

\subsection{Accountability and Presentation Results}

A low-count contacts record without processing context returns insufficient evidence rather than no concern. Its communication priority is standard because the absence of legal context warrants review. A low-count location record with purpose, contract basis, controller identity, one-day retention, and user-initiated context returns no technical concern and a silent priority.

A microphone record marked as consent-based, special-category, and accessed after consent withdrawal returns the urgent review status. The fixture does not prove a real GDPR infringement; it confirms that the pack's stated escalation branch is implemented. A special-category microphone record without an Article 9 condition lists that condition among missing evidence. An invented lawful basis and invalid retention value are rejected at admission.

Dashboard projection preserves three different states from the accountability fixtures. Summary totals show one action item, one evidence gap, and one no-concern result. Filters surface the expected finding. Every presentation state has both a textual label and marker, ensuring that source-level semantics do not depend on colour alone.

\subsection{Regulation-Pack Results}

The registry exposes at least two packs and recognises only the two expected identifiers. A forged identifier is rejected. The same location audit with a count of 40 is evaluated under both packs. The resulting findings retain different regulation identifiers and different thresholds because each pack supplies its own baseline.

This test is stronger than checking that a dropdown changes text. It passes the same admitted evidence through two independently instantiated engines and inspects the resulting finding objects. The Research Baseline result also contains a caveat stating that it is not law. The test does not validate either pack's legal interpretation; it validates decoupling, identity, and disclosure.

Future pack validation should add a contract suite automatically applied to every registry entry. It should require a unique identifier, version, official-source policy, complete supported-rule set, positive finite baselines and multipliers, non-empty rationale and caveat, deterministic missing-evidence logic, and fixtures reviewed by the responsible domain expert.

\subsection{Build Reproduction}

The release build initially encountered Windows path and dependency-layout problems. A hoisted local dependency layout allowed native CMake paths to remain within workable bounds. React Native's new architecture was disabled for the release candidate because enabling it in this Windows checkout reproduced path-related native build failures. This is a build-environment decision, not evidence that the legacy architecture is preferable for future releases.

Gradle reported deprecated features that will require attention before Gradle 9 compatibility. It also reported an SDK XML version warning caused by differences between installed Android Studio and command-line components. Neither warning prevented the recorded release build, but both are retained as maintenance debt rather than suppressed.

The release task generated bundles for multiple Android ABIs. The APK was installed on the ARM64 physical device. The AAB was not uploaded to Play Console, processed by Play signing, or delivered through an internal track. Its success establishes local bundle generation only.

Hashes were recorded immediately after the final permission-manifest rebuild. Any source or build-tool change requires rebuilding and computing new hashes. The files are generated artefacts and are not treated as immutable simply because their names remain the same.

\subsection{Manifest and Installed-Package Reconciliation}

Source-manifest inspection alone was insufficient because dependencies contribute declarations through manifest merging. The initial merged manifest exposed legacy external-storage permissions inherited from transitive Expo components. Revision 10 added explicit merger removals for read and write external storage, rebuilt both artefacts, and inspected the merged result again.

The final merged manifest contained Internet, wake-lock, network-state, boot-completed, foreground-service, and an application-scoped dynamic-receiver signature permission. These operational declarations are consistent with linking and WorkManager components. The installed-package dump reported the same requested set and no runtime permission grants. This two-level check supports the documentation statement that no sensitive runtime permission is requested.

The package dump also confirmed application identity, version name, version code, minimum SDK, target SDK, primary ABI, install state, and signing version. The release used the QA fallback signing identity because production environment variables were absent. A successful package inspection is therefore not proof of final upload-key configuration.

\subsection{Physical Interaction Protocol}

The device was identified by ADB before installation. The final release APK was streamed with replacement enabled, the package was force-stopped, and the launcher intent was injected. Package metadata was read after installation so that a stale or differently named package could not be mistaken for the new build.

Interaction evidence included screenshots and UI hierarchy dumps for overview, findings, and settings. A controlled demonstration populated the dashboard and finding list. The final overview was captured again after the permission-manifest rebuild and installation. Runtime review used filtered AndroidRuntime and ReactNativeJS error channels; no matching fatal output appeared in the bounded recent log window.

An important measurement issue arose from coordinate spaces. The physical display reported 1080 by 2400 pixels, while the captured application window and hierarchy used 1080 by 2244. Initial taps derived from the full display therefore appeared to miss navigation targets. A bounded coordinate sweep and hierarchy inspection established the offset, after which navigation succeeded. The event is reported because excluding it would conceal a real threat to ADB-based UI-test validity.

Future automation should select accessibility nodes or stable resource identifiers instead of relying on raw coordinates. It should record display size, application bounds, system-bar insets, orientation, density, and font scale for every run. Screenshots should be associated with the installed package hash and build hash so that visual evidence cannot drift from the tested binary.

\subsection{Visual Review Method}

Visual assessment considered hierarchy, identity, first-action clarity, state distinction, provenance visibility, limitation placement, empty states, populated states, navigation consistency, destructive-action treatment, and overall coherence. The assessment did not score aesthetic preference numerically because no validated participant scale was administered.

The overview succeeded by making the product proposition and review action dominant while keeping the evidence-reach warning in the first viewport. Summary cards used count, icon, label, and colour. Findings showed synthetic source and active pack near the decision. Settings gave the pack caveat and controlled demonstration their own sections. The custom navigation remained legible against the off-white background and used a clear selected state.

The result satisfies a project-defined candidate gate rather than a universal usability standard. One evaluator and one physical device cannot establish willingness to adopt across a population. The strongest supported statement is that the rendered candidate is coherent, presents its boundaries, and is suitable for internal-track and participant evaluation.

\subsection{Result Boundary}

The evidence supports the claim that the specified code compiles, the rule engine conforms to its fixtures, the Android artefacts build, the final package installs and launches, the main interface renders coherently, and provenance remains visible in the tested workflows. It does not support claims about 24-hour execution, battery cost, memory leaks, unrestricted third-party observation, multiple OEMs, legal compliance, accessibility conformance, production signing, or store acceptance.

\subsection{Evidence Triangulation}

No single method in the campaign was designed to carry the complete acceptance claim. Source inspection is good at confirming constants, branches, labels, and data flow, but it cannot establish which code entered an APK or how a screen rendered. Automated tests execute specified behaviours repeatedly, but their assertions can share assumptions with the implementation. Package inspection observes the assembled artefact, but it cannot prove the meaning of every runtime path. Device interaction observes a concrete installation, but the sample is narrow. The evaluation therefore uses convergence among methods while preserving the scope of each.

The deterministic tests and source review converge on threshold semantics. Source review identifies the pack-provided baseline, multiplier, daily threshold, burst limit, and source restrictions. Boundary assertions then test values immediately around those decision points. The fixed corpus tests varied inputs at scale. Agreement across these methods reduces the chance that the reported rule differs accidentally from the implemented one. It still does not establish that the rule is a valid legal classifier.

Runtime-boundary tests and presentation tests converge on uncertainty handling. Malformed evidence is rejected at admission, low-context evidence becomes an evidence gap, and presentation projection retains that state. This end-to-end check matters because a correct engine could be undermined if the UI translated an unknown state into reassuring copy. The tests show that selected fixtures preserve meaning through the current projection; screenshot review confirms the principal labels are visible in the rendered flow.

Build logs, output existence, checksums, manifest inspection, installed-package metadata, and physical launch converge on artefact identity. Each closes a different substitution risk. A successful Gradle message without an output file would be insufficient. A local file without a checksum could be confused after rebuilding. Source manifest facts could differ after dependency merging. An installed package could be a stale version. Reading these signals together supports the conclusion that the evaluated handset ran the recorded release candidate.

Visual review and UI-hierarchy inspection also complement each other. A screenshot shows layout, colour, clipping, and hierarchy as a person sees them. A hierarchy dump exposes text, bounds, and accessibility-node structure that may not be obvious visually. Neither replaces screen-reader interaction, large-font testing, or participant comprehension. Their combined use supports a bounded expert assessment and identifies the next accessibility evidence required.

Triangulation is not vote counting. If one authoritative check contradicts another, the discrepancy must be resolved rather than outvoted. For example, an installed-package dump showing a sensitive permission would invalidate a source-only conclusion even if several documents said otherwise. The evaluation gives greater weight to evidence closest to the released artefact for package claims and to independent human or legal review for substantive claims.

\subsection{Oracle Independence and Correlated Defect Risk}

The main corpus labels are derived from the same published threshold model the implementation is intended to realise. This is appropriate for conformance testing: it asks whether the program implements the specification. It is insufficient for empirical validation because a mistaken specification can yield perfect agreement. The measured 1.0000 precision and recall are therefore implementation-oracle agreement values, not estimates of performance on real legal cases.

Several design choices reduce accidental coupling. The generator records expected polarity directly from its construction, a fresh engine limits cross-case state, boundary cases are written explicitly, and extended tests target properties beyond the main daily-total rule. Nevertheless, developer knowledge shaped both sides. The campaign did not use a separately developed executable oracle, blinded labels, or an external benchmark.

Mutation testing would provide one useful next check. Deliberately changing a comparison from greater-than to greater-than-or-equal, weakening source validation, removing replay replacement, or allowing overlapping history should cause specific tests to fail. Surviving mutations would identify assertions that execute code without distinguishing correct behaviour. Property-based generation could then explore many valid and invalid shapes while preserving stated invariants.

An independent reference implementation could further reduce correlated defects. It should be written from the regulation-pack contract by a reviewer who has not read the production engine, ideally in a different language or style. Both implementations could evaluate a versioned interchange corpus. Disagreement would trigger case analysis rather than majority voting. This process would validate formal semantics, though it would still not validate legal relevance.

For legal or ecological validation, the oracle must change. Controlled event traces can provide ground truth about whether acquisition observed an action. Independent privacy reviewers can label whether a scenario warrants investigation, using a documented protocol and recording disagreement. Neither source should be described as perfect legal truth. The evaluation question, sampling frame, and annotation uncertainty must be reported alongside performance metrics.

This distinction explains why the thesis retains the term deterministic corpus rather than benchmark. A benchmark normally suggests comparison against an external reference population. The current corpus is a rigorous regression artefact that protects rule semantics as the code changes. Its value is high, but different from an independently labelled real-world set.

\subsection{Metric Interpretation and Error Analysis}

Precision and recall were included because they expose false-positive and false-negative cells more clearly than accuracy in an imbalanced corpus. The designed 80/20 mixture would allow a trivial always-positive strategy to appear 80 per cent accurate. The full confusion matrix prevents that reading: all 200 controls were negative and all 800 constructed signals were positive under the specified oracle.

The exact totals also make the experiment reproducible. The campaign did not sample a changing external population; it evaluated a fixed generated set. Confidence intervals around the observed proportions would describe repeated sampling from an assumed generator, not uncertainty about the fixed run. More importantly, they would not address specification validity. For this reason the report presents exact agreement and claim boundaries instead of using a narrow interval to imply real-world certainty.

Error analysis remains meaningful even when the main matrix has no disagreements. Adversarial suites reveal categories the primary generator does not represent: unsafe numeric values, prototype-derived identifiers, malformed context, source forgery, temporal overlap, replay, and history isolation. These tests show that corpus size alone is not coverage. A small, carefully targeted case can provide more security evidence than hundreds of additional values drawn from the same safe range.

Future field studies should report more than aggregate precision and recall. Results should be stratified by permission, application category, source authority, foreground status, OEM, Android version, and window resolution. Reviewers should inspect false positives and negatives for common causes. Missing acquisition records should not be treated as model negatives. Cases excluded because evidence is insufficient should be counted and explained rather than removed from the denominator without notice.

Calibration should be evaluated separately from discrimination. A status hierarchy may rank more concerning cases correctly while its displayed wording still conveys excessive certainty. Participant studies can test whether confidence changes appropriately between authoritative, synthetic, and incomplete evidence. This human calibration outcome is central to the product's trust claim and cannot be inferred from the engine matrix.

Operational metrics also need denominators. A launch success rate requires a defined number of clean starts and devices. A crash-free percentage requires a meaningful observation period and collection method. A jank percentage depends on a representative interaction trace and enough rendered frames. The current evaluation reports bounded events directly rather than producing population-style percentages from a short sample.

\subsection{Store-Candidate Evidence Review}

The project-defined store-candidate gate was assessed across product identity, core journey, disclosure, packaging, privacy posture, accessibility foundations, and operational ownership. Product identity was supported by a consistent name, icon, visual language, and package identifier. Core journey was supported by first-launch overview, findings navigation, detailed rationale, pack selection, demonstration control, and local-data deletion. Disclosure was supported by visible source labels, pack identity, evidence limits, and non-verdict language.

Packaging evidence included release APK and AAB generation, target SDK 36, ABI outputs, final checksums, manifest reconciliation, and installation of the APK. Privacy posture included no sensitive runtime permission in the final package, disabled backup, local-first storage, and a draft privacy policy consistent with the tested build. Accessibility foundations included textual state labels, non-colour cues, labelled controls, target sizing intent, and a simple navigation structure.

The gate was not interpreted as proof that every store prerequisite is complete. Production upload signing was absent. The privacy policy required owner-controlled HTTPS publication. Console declarations, listing assets, support details, content rating, internal-track delivery, and store review remained external. These gaps do not negate the engineering candidate; they determine why it must not be described as a publicly released production version.

The visual adjectives in the gate were operationalised cautiously. ``Beautiful'' meant coherent hierarchy, spacing, typography, colour restraint, branded identity, and absence of visible broken states in the inspected flows. ``Trustworthy'' meant that evidence provenance and limitations were prominent, destructive controls were explicit, and the product did not present a legal badge. ``Willing to use'' meant that an evaluator could understand the purpose, reach findings, inspect reasons, switch the review framework, and control data without encountering a blocking usability defect.

These definitions create a repeatable expert gate, but they do not convert subjective judgement into population evidence. A single evaluator can miss cultural preferences, accessibility barriers, unfamiliar terminology, or long-term fatigue. The appropriate next step is internal-track and participant evaluation using the actual distributed candidate. The expert gate answers whether that investment is justified, not whether the market has accepted the product.

The candidate judgement also depends on the declared product category. A local research and accountability aid can be useful with bounded imported or simulated evidence. A product marketed as universal device surveillance would fail because the acquisition authority is not established. Store copy must preserve the former positioning. Changes to marketing are therefore capable of invalidating readiness even if the binary remains identical.

\subsection{Evaluation Repeatability and Change Control}

The evidence applies to a specific source revision and generated artefacts. Repeating only the final UI flow after changing engine code would not renew rule evidence. Rebuilding only the APK after changing privacy copy would not renew screenshot evidence. The project therefore treats acceptance as a dependency graph: a change invalidates every downstream artefact that incorporates or presents it.

Engine or pack changes require TypeScript compilation, logical suites, boundary suites, presentation projections, release rebuilds, hashes, and targeted device review. Android manifest or Gradle changes require merged-manifest inspection, package dump reconciliation, rebuilds, installation, and launch. UI changes require lint, package build, hierarchy and screenshot capture, accessibility checks, and renewed expert review. Policy or store-description changes require reconciliation with package behaviour and data flows.

Reproducible scripts reduce omission but do not remove the need to inspect their outputs. A command may exit successfully while testing stale generated code, selecting the wrong device, or reading a prior file. The campaign therefore compiles generated runners immediately before execution, checks package identity after installation, and associates evidence with checksums. Future automation should emit a machine-readable manifest linking commit, tool versions, commands, artefact hashes, device facts, and screenshots.

Environment details also matter. Windows path length and native dependency layout affected release compilation. The chosen workaround and architecture flag belong in the build record because a different environment may produce different outputs. Tool warnings about future Gradle compatibility remain maintenance observations even when the current build passes.

A release tag should be immutable, while packs and policies use explicit version identifiers. If a later investigation finds an error, the project should publish a successor and a correction note rather than rewriting the historical evidence. This thesis follows the same principle by creating Revision 10 chapter files instead of replacing Revision 9. Preservation makes it possible to see which claims changed and why.

Repeatability finally includes honest failure recording. Network, dependency, device-coordinate, and build-path failures can be part of the evidence because they expose operational assumptions. Removing them from the narrative would make future reproduction harder. A trustworthy evaluation explains both the final pass and the conditions needed to obtain it.
